%%%%%%%%%%%%%%%%%%%%%%%%%%%%%%%%%
%
%       EXERCÍCIO
%
%%%%%%%%%%%%%%%%%%%%%%%%%%%%%%%%%

\definevalorquestao

\begin{exercicioBanco}[\valorquestao]
Uma população de bactérias em um experimento cresce de acordo com a função \(N(t) = 125 \cdot 5^{0,6t}\), em que \(N(t)\) representa a quantidade de bactérias e \(t\) é o tempo decorrido em horas. O tempo necessário para que essa população atinja \(\n{15625} = 5^{6}\) bactérias é de:

\newcommand{\alternativas}{%
    \begin{center}
        \begin{tabularx}{\textwidth}{XXXXX}
            (a) \(3\) horas. &
            (b) \(4\) horas. &
            (c) \(5\) horas. &
            (d) \(6\) horas. &
            (e) \(8\) horas.
        \end{tabularx}
    \end{center}
}

\newcommand{\respostaUm}{C}

% Lógica condicional para exibição
\ifdefstring{\atividade}{lista}{%
    \alternativas
    \vspace{0.5em}
    
    \noindent\textbf{Resposta:} letra \textbf{\resposta}.
}{%
    \ifdefstring{\modo}{objetivo}{%
        \alternativas
    }{%
        % modo = discursiva → não mostra alternativas
    }
}
\end{exercicioBanco}

